% \documentclass[a4paper,12pt]{article}
% \usepackage[ngerman]{babel}
% \usepackage[utf8]{inputenc}
% \usepackage{amsmath,amssymb}
% \usepackage{geometry}
% \geometry{a4paper, margin=2.5cm}
% \usepackage{enumitem}
% \usepackage{hyperref}
% \usepackage{fancyhdr}
% \pagestyle{fancy}
% \fancyhead[L]{Mathematik für Ingenieure}
% \fancyhead[C]{Arbeitsblatt: Produktregel}
% \fancyhead[R]{\today}

% \begin{document}

\begin{center}
    {\textbf{Arbeitsblatt: Die Produktregel beim Ableiten}} \\[1em]
    {Mathematik für Ingenieure}
\end{center}

\section*{1. Einleitung und Motivation}
In vielen ingenieurwissenschaftlichen Anwendungen treten Funktionen auf, die als Produkt zweier (oder mehrerer) Funktionen dargestellt werden können. Die Produktregel ist ein zentrales Werkzeug, um solche Funktionen korrekt abzuleiten.

\section*{2. Mathematische Herleitung der Produktregel}
Seien $f(x)$ und $g(x)$ differenzierbare Funktionen. Die Produktregel lautet:
\begin{align*}
    (f \cdot g)'(x) = f'(x) \cdot g(x) + f(x) \cdot g'(x)
\end{align*}
\textbf{Herleitung:}
\begin{align*}
    (f \cdot g)'(x) &= \lim_{h \to 0} \frac{f(x+h)g(x+h) - f(x)g(x)}{h} \\
    &= \lim_{h \to 0} \frac{f(x+h)g(x+h) - f(x)g(x+h) + f(x)g(x+h) - f(x)g(x)}{h} \\
    &= \lim_{h \to 0} \left[ \frac{f(x+h)-f(x)}{h}g(x+h) + f(x)\frac{g(x+h)-g(x)}{h} \right] \\
    &= f'(x)g(x) + f(x)g'(x)
\end{align*}

\section*{3. Beispiele}
\subsection*{Beispiel 1}
Gegeben sei $f(x) = x^2$ und $g(x) = \sin(x)$. Berechnen Sie die Ableitung von $h(x) = f(x) \cdot g(x) = x^2 \sin(x)$.
\\
\textbf{Lösung:}
\begin{align*}
    h'(x) &= (x^2)' \cdot \sin(x) + x^2 \cdot (\sin(x))' \\
    &= 2x \sin(x) + x^2 \cos(x)
\end{align*}

\subsection*{Beispiel 2}
$h(x) = e^x \cdot \ln(x)$
\\
\textbf{Lösung:}
\begin{align*}
    h'(x) &= (e^x)' \cdot \ln(x) + e^x \cdot (\ln(x))' \\
    &= e^x \ln(x) + e^x \cdot \frac{1}{x} \\
    &= e^x \left( \ln(x) + \frac{1}{x} \right)
\end{align*}

\section*{4. Typische Fehlerquellen}
\begin{itemize}
    \item Die Produktregel wird mit der Kettenregel verwechselt.
    \item Es wird nur ein Summand der Produktregel verwendet.
    \item Die Ableitungen der Einzelterme werden falsch berechnet.
\end{itemize}

\section*{5. Anwendungsaufgaben}
\subsection*{Aufgabe 1}
Leiten Sie $h(x) = (3x^2 - 2x) \cdot \cos(x)$ ab.

\subsection*{Aufgabe 2}
Berechnen Sie die Ableitung von $h(x) = x \cdot e^{2x}$.

\subsection*{Aufgabe 3}
Gegeben ist $h(x) = (x^3 + 1) \cdot \ln(x)$.

\subsection*{Aufgabe 4}
Bestimmen Sie die Ableitung von $h(x) = \sqrt{x} \cdot \arctan(x)$.

\subsection*{Aufgabe 5}
$h(x) = (x^2 + 2x) \cdot (x^3 - x)$

\section*{6. Lösungen zu den Anwendungsaufgaben}
\begin{enumerate}[label=\textbf{Aufgabe \arabic*:}]
    \item $h'(x) = (6x - 2)\cos(x) + (3x^2 - 2x)(-\sin(x))$
    \item $h'(x) = 1 \cdot e^{2x} + x \cdot 2e^{2x} = e^{2x} + 2x e^{2x}$
    \item $h'(x) = 3x^2 \ln(x) + (x^3 + 1) \cdot \frac{1}{x}$
    \item $h'(x) = \frac{1}{2\sqrt{x}} \arctan(x) + \sqrt{x} \cdot \frac{1}{1+x^2}$
    \item $h'(x) = (2x+2)(x^3-x) + (x^2+2x)(3x^2-1)$
\end{enumerate}

\section*{7. Aufgaben zur Selbstkontrolle}
\begin{itemize}
    \item $h(x) = x^2 \cdot e^x$
    \item $h(x) = \sin(x) \cdot \ln(x)$
    \item $h(x) = x \cdot \arcsin(x)$
    \item $h(x) = (x^2+1) \cdot \cos(x)$
    \item $h(x) = e^x \cdot \arctan(x)$
\end{itemize}

\vfill
\textbf{Weiterführende Literatur:} Siehe z.B. Bronstein, Taschenbuch der Mathematik oder das Skriptum zur Mathematik für Ingenieure.

